\subsection{Modeling phylodynamic rates with a neural network}

Phylodynamic inference seeks to quantify the population dynamic and sampling processes that give rise to an observed phylogenetic tree $T$; we refer to the model of this process as the \emph{phylodynamic model}.
In macroevolutionary applications, the relevant parameters typically include speciation, extinction, and fossilization rates, whereas in epidemiological settings they may include transmission, recovery, and sampling rates.
A widely used model for population dynamics is the birth--death model, parameterized by a birth rate $\lambda$, a death rate $\mu$, and a sampling rate $\psi$ \cite{stadler2013uncovering}.
More elaborate birth--death models may additionally incorporate migration processes and type-dependent rates \cite{Vaughan2025BDMM}.
While this work focuses on phylodynamic inference under the birth--death framework, the concepts developed here are also applicable to the coalescent framework \cite{kingman1982coalescent}.

The posterior distribution over the phylodynamic parameters is given by
\[
    P(\theta \mid T)
    \propto
    P(T \mid \theta)\,
    P(\theta),
\]
where $\theta = (\lambda, \mu, \ldots)$ denotes the set of parameters of the birth--death model.
The phylodynamic likelihood $P(T \mid \theta)$ specifies the probability of observing the tree topology and branch lengths under the assumed birth--death process.
Phylodynamic inference therefore aims to identify the parameter values that most plausibly generated a given phylogenetic tree, while naturally accounting for uncertainty through the posterior distribution.
In practice, phylogenetic trees are typically not observed directly but are themselves inferred from empirical data, such as aligned molecular sequences or morphological characters.
Consequently, the phylodynamic process is commonly used as a tree prior within a joint Bayesian framework in which the phylogeny is estimated simultaneously with the phylodynamic parameters.

Most phylodynamic models directly estimate the parameters of a birth--death process, such as speciation and extinction rates in macroevolution, or migration rates and reproduction numbers in epidemiology.
Alternative approaches instead model these rates as functions of external predictors.
In GLM–based formulations, phylodynamic parameters are given by a (tipically logarithmic) link function applied to a linear combination of predictors.
The inferred parameters are the intercept and regression coefficients of this relationship \cite{lemey2014unifying, muller2019inferring, valenzuela2021comprehensive}.
Here, we model variation in phylodynamic rates as a function of predictors using a feedforward neural network.

A feedforward neural network (FFNN) is a flexible function approximator capable of capturing complex, nonlinear relationships between inputs and outputs \cite{haykin1994neural}.
An FFNN consists of an input layer, one or more hidden layers, and an output layer.
Each layer applies a linear transformation followed by a nonlinear activation function.
Formally, for layer $l$ with input vector $\mathbf{x}^{(l-1)}$, weight matrix $\mathbf{W}^{(l)}$, and bias vector $\mathbf{b}^{(l)}$, the output is given by
\[
    \mathbf{x}^{(l)} = f\!\left(\mathbf{W}^{(l)} \mathbf{x}^{(l-1)} + \mathbf{b}^{(l)}\right),
\]
where $f(\cdot)$ denotes a nonlinear activation function, such as the rectified linear unit (ReLU) or the sigmoid function.
By stacking multiple layers, an FFNN can approximate arbitrary continuous functions, enabling it to learn complex mappings from predictor variables to phylodynamic rates.

In standard applications, FFNNs are used in supervised learning settings, where the network is trained to map inputs to observed target values by minimizing a loss function, such as mean squared error for regression or cross-entropy for classification.
In this context, the network weights and biases are treated as parameters that are optimized during training.
In contrast, BELLA adopts an unsupervised use of FFNNs to integrate external, non-phylogenetic data into phylodynamic analyses.
The FFNN maps these input predictors to a subset of the phylodynamic parameters, denoted \( \theta_{\text{FFNN}} \), while priors are placed on the network weights \( \mathbf{W} \) and biases \( \mathbf{b} \).
The remaining parameters, \( \theta_{\text{base}} \), are not predicted by the FFNN and are instead sampled directly from their standard priors.
This separation allows the model to exploit external covariates for parameters that can be informed by additional data, while retaining classical Bayesian sampling for the remaining parameters.

To ensure that the FFNN produces parameter values within biologically or mathematically meaningful ranges, the choice of output activation functions can be tailored to the specific parameter being modeled.
For instance, parameters that are inherently bounded between zero and one, such as probabilities, can be mapped through a sigmoid transformation.
More generally, logistic functions with user-specified upper and lower bounds allow parameters to be constrained to finite intervals, while exponential or \texttt{softplus} activations can enforce strict positivity.
By aligning the output activation function with the natural support of each parameter, the model avoids generating implausible values and improves the efficiency of posterior exploration.

As in standard phylodynamic analyses, Markov chain Monte Carlo (MCMC)~\cite{CHIB20013569} is used to jointly sample from the posterior
\[
    P(\mathbf{W}, \mathbf{b}, \theta_{\text{base}} \mid T)
    \propto
    P(T \mid \theta_{\text{FFNN}}(\mathbf{W}, \mathbf{b}), \theta_{\text{base}}) \,
    P(\mathbf{W}, \mathbf{b}) \,
    P(\theta_{\text{base}}),
\]
with operators acting on both the network parameters and the standard prior-sampled parameters.
At each MCMC iteration, new values for $\mathbf{W}$ and $\mathbf{b}$ are proposed and used by BELLA to map the input predictors to the phylodynamic rates $\theta_{\text{FFNN}}$, which are then combined with $\theta_{\text{base}}$ to evaluate the phylodynamic likelihood.
When phylodynamic rates vary across lineages, the predictor values associated with individual branches are typically unknown.
For instance, the geographic location or trait state of an ancestral branch is rarely observed.
In such cases, the FFNN maps predictor values to lineage-specific rates for each possible ancestral state, and the likelihood is computed by integrating over all these possibilities \cite{Vaughan2025BDMM}.
This procedure naturally propagates uncertainty about ancestral predictor values through the neural network and into the posterior distribution, providing a fully Bayesian integration of external data into phylodynamic inference.
